% Title page for double-anonymous review.
% Kept as a SEPARATE file from sentinel_e.tex (the blinded manuscript) so the
% manuscript file itself carries no author-identifying information. Submit
% both files together; do not merge this content back into sentinel_e.tex.
\documentclass[11pt,a4paper]{article}

\usepackage[T1]{fontenc}
\usepackage[utf8]{inputenc}
\usepackage{lmodern}
\usepackage[margin=2.5cm]{geometry}
\usepackage{amsmath,amssymb}
\usepackage{microtype}
\PassOptionsToPackage{hidelinks}{hyperref}
\usepackage{hyperref}
\usepackage{orcidlink}

\graphicspath{{../results/figures/}}
% Auto-generated by experiments/make_macros.py -- do not edit.
% Every number quoted in the manuscript prose is defined here and comes
% directly from the JSON payloads written by the experiment scripts.

\newcommand{\ExpOneHorizonReps}{200}
\newcommand{\ExpOneReps}{1500}
\newcommand{\HorizonFixedPfa}{1.000}
\newcommand{\HorizonFixedPfaPct}{100}
\newcommand{\HorizonFixedPfaShort}{0.640}
\newcommand{\HorizonMaxFrames}{240{,}000}
\newcommand{\HorizonMaxHours}{2.7}
\newcommand{\HorizonMinHours}{0.2}
\newcommand{\HorizonPvaluePfa}{1.000}
\newcommand{\HorizonPvaluePfaPct}{100}
\newcommand{\HorizonSentinelPfa}{0.000}
\newcommand{\ValidityAlphaLoose}{0.20}
\newcommand{\ValidityMaxSentinelPfa}{0.001}
\newcommand{\ValidityPfaOracleLoose}{0.044}
\newcommand{\ValidityPfaPvalue}{0.998}
\newcommand{\ValidityPfaPvaluePct}{100}
\newcommand{\ValidityPfaSentinel}{0.000}
\newcommand{\ValidityPvalueArl}{11{,}155}
\newcommand{\ValidityPvalueArlMin}{7.4}
\newcommand{\ValidityTightestAlpha}{0.005}
\newcommand{\ValidityTightestCiHi}{0.003}
\newcommand{\ValidityTightestPfa}{0.000}
\newcommand{\ValidityWorstCiHi}{0.004}

% Placeholders for experiments that have not been run yet.  These make
% the manuscript compile from a partial results directory; a '??' in the
% PDF means the corresponding experiment output is missing.
\providecommand{\AblChangepointDelay}{\textbf{??}}
\providecommand{\AblChangepointMiss}{\textbf{??}}
\providecommand{\AblChangepointPfa}{\textbf{??}}
\providecommand{\AblDilationDelay}{\textbf{??}}
\providecommand{\AblDkwDelay}{\textbf{??}}
\providecommand{\AblDkwSpeedup}{\textbf{??}}
\providecommand{\AblEveryFramePfa}{\textbf{??}}
\providecommand{\AblFullDelay}{\textbf{??}}
\providecommand{\AblFullMiss}{\textbf{??}}
\providecommand{\AblLag}{\textbf{??}}
\providecommand{\AblNoCcPfa}{\textbf{??}}
\providecommand{\AblNumInvalid}{\textbf{??}}
\providecommand{\AblNumVariants}{\textbf{??}}
\providecommand{\AblPointPriorMiss}{\textbf{??}}
\providecommand{\AblPooledPfa}{\textbf{??}}
\providecommand{\AblRawCalPfa}{\textbf{??}}
\providecommand{\CostBackboneMs}{\textbf{??}}
\providecommand{\CostBackboneParams}{\textbf{??}}
\providecommand{\CostFitSec}{\textbf{??}}
\providecommand{\CostRealtimeFactor}{\textbf{??}}
\providecommand{\CostSentinelPct}{\textbf{??}}
\providecommand{\CostSentinelUs}{\textbf{??}}
\providecommand{\CoverDaytimeAbstain}{\textbf{??}}
\providecommand{\CoverDaytimeMiss}{\textbf{??}}
\providecommand{\DelayInadmissibleList}{\textbf{??}}
\providecommand{\DelayNumInadmissible}{\textbf{??}}
\providecommand{\DelaySentinelCi}{\textbf{??}}
\providecommand{\DelaySentinelMiss}{\textbf{??}}
\providecommand{\DelaySentinelPfa}{\textbf{??}}
\providecommand{\DelaySentinelSec}{\textbf{??}}
\providecommand{\DiagAcfRaw}{\textbf{??}}
\providecommand{\DiagAcfThin}{\textbf{??}}
\providecommand{\DiagLjungRaw}{\textbf{??}}
\providecommand{\DiagLjungThin}{\textbf{??}}
\providecommand{\DiagPmin}{\textbf{??}}
\providecommand{\EpiHardPiCpMiss}{\textbf{??}}
\providecommand{\EpiHardPiEpiMiss}{\textbf{??}}
\providecommand{\EpiManyCpDelay}{\textbf{??}}
\providecommand{\EpiManyCpMiss}{\textbf{??}}
\providecommand{\EpiManyEpiDelay}{\textbf{??}}
\providecommand{\EpiManyEpiMiss}{\textbf{??}}
\providecommand{\EpiMaxCount}{\textbf{??}}
\providecommand{\EpiMinPi}{\textbf{??}}
\providecommand{\EpiOneCpDelay}{\textbf{??}}
\providecommand{\EpiOneCpMiss}{\textbf{??}}
\providecommand{\EpiOneEpiDelay}{\textbf{??}}
\providecommand{\EpiOneEpiMiss}{\textbf{??}}
\providecommand{\ExpTwoReps}{\textbf{??}}
\providecommand{\NetCameras}{\textbf{??}}
\providecommand{\NetFdrLevel}{\textbf{??}}
\providecommand{\NetFleets}{\textbf{??}}
\providecommand{\NetGnnDelayReduction}{\textbf{??}}
\providecommand{\NetGnnDelaySec}{\textbf{??}}
\providecommand{\NetGnnFdr}{\textbf{??}}
\providecommand{\NetGnnNullPfa}{\textbf{??}}
\providecommand{\NetGnnPower}{\textbf{??}}
\providecommand{\NetNoGraphDelaySec}{\textbf{??}}
\providecommand{\NetNoGraphNullPfa}{\textbf{??}}
\providecommand{\ShiftMaxRatio}{\textbf{??}}
\providecommand{\ShiftMondrianPfa}{\textbf{??}}
\providecommand{\ShiftPooledPfa}{\textbf{??}}
\providecommand{\ShiftResidualPfa}{\textbf{??}}


\title{\textbf{SENTINEL-E}: Anytime-Valid Sequential Detection for\\
       Streaming Visual Anomaly Surveillance}

\author{%
  Quang-Vinh Dang\,\orcidlink{0000-0002-3877-8024}\\[2pt]
  \small British University Vietnam, Hung Yen, Vietnam\\
  \small \texttt{vinh.dq4@buv.edu.vn}
}
\date{}

\begin{document}
\maketitle
\thispagestyle{empty}

\begin{center}
  \small\emph{Title page for double-anonymous review. The accompanying
  manuscript file (\texttt{sentinel\_e.tex}/\texttt{.pdf}) carries no author
  information; this page and its declarations are for editorial use and are
  withheld from reviewers. Competing interests, data/code availability, and
  ethical approval are also stated, in non-identifying form, in the
  manuscript's own Declarations section; only funding and author
  contributions appear on this page alone.}
\end{center}
\vspace{1em}

\begin{abstract}
\noindent
Streaming visual-anomaly systems --- illegal dumping among them --- are
trained and evaluated as fixed-sample-size classifiers: given a short clip,
emit one decision and report clip-level precision and recall. Deployment does
not look like that: a municipal camera runs for weeks, and what an operator
manages is the number of false alarms accumulated over the whole monitoring
period, because each one costs a dispatch. Alerting the first time a
per-frame threshold is crossed on a continuous stream is optional stopping,
and it inflates the true false-alarm rate by an amount clip-level evaluation
cannot show: a per-frame rule fed p-values individually valid at the $1\%$
level raises at least one false alarm on \ValidityPfaPvaluePct\% of
$40$-minute null streams.
%
Our contribution is an algorithm. Every anytime-valid sequential change
detector in use --- CUSUM, Shiryaev--Roberts, e-detectors, non-partitioned
e-detectors --- models a switch that, once it happens, lasts forever. A
dumping event is an \emph{episode}: it begins, flickers under intermittent
occlusion, ends, and recurs weeks later at the same site. We replace the
mixture over changepoint locations with a mixture over episodes, realised as
a two-state Markov chain whose forward recursion collapses a sum over $2^t$
trajectories into two numbers. The resulting \emph{episodic e-process} is an
exact non-negative supermartingale --- so Ville's inequality gives
time-uniform false-alarm control with no new machinery --- costs $O(1)$ per
frame against $O(t)$ for a general mixture over active starts, and contains
the changepoint mixture exactly as the boundary case in which episodes never
end, a strict generalisation along an axis the anytime-valid e-process
family has left fixed at infinity. Around this core sits a per-camera
conformal layer calibrated conditionally on observable context and corrected
by exact Beta order-statistic levels, and a fleet layer in which a graph
network steers each camera's episode hazard through a \emph{predictable}
modulation --- provably unable to damage the guarantee --- combined across
cameras by e-BH under arbitrary dependence.
%
Two evidence sources back these claims. Large-horizon simulated streams
($10^5$--$10^6$ frames) show every classical and parametric baseline we test
--- fixed threshold, per-frame p-value threshold, CUSUM, Shiryaev--Roberts, a
parametric e-detector, E-SHIFT --- failing to hold its nominal false-alarm
level, while SENTINEL-E and the changepoint mixture it strictly contains both
hold it; a paired comparison on a shared betting grid resolves a real
advantage for the episodic mixture at two and four recurring episodes and a
real, smaller cost at one. We also report where the guarantee stops holding:
an unmodelled covariate shift beyond $\times\ShiftResidualBreaks$ inflates the
rate unless a shift-weighting safeguard is engaged, at a steep cost in power.
%
New in this revision: the pipeline is also run, unmodified, on real
surveillance video (CUHK Avenue and UCSD Ped2) with a real optical-flow
backbone rather than a simulator. This surfaces a limitation the simulated
streams could not: real per-frame scores are autocorrelated across
hundreds of frames (\RealAvenueFullLag{} on Avenue, \RealPedTwoFullLag{} on
Ped2), close to or beyond a real episode's own duration, leaving only a
narrow lag window in which validity and
power coexist at all. Inside that window the guarantee holds on both real
corpora --- realised false-alarm rate \RealPedTwoPfaHeadline{} on Ped2 and
\RealAvenuePfaHeadline{} on Avenue at $\alpha=0.1$ --- while every classical
baseline tested fails to hold it on the same real streams; detection power is
real but modest on Ped2 and weak on Avenue, whose real episodes are shorter.
This gap, not the simulated results, is the paper's central open problem.
SENTINEL-E costs \CostSentinelUs\,\textmu s per frame, \CostSentinelPct\% of
the detector it wraps.
\end{abstract}

\vspace{0.5em}
\noindent\textbf{Keywords:} video anomaly detection; illegal dumping detection;
anytime-valid inference; e-values; conformal prediction; sequential change
detection; false discovery rate; video surveillance; edge deployment

\section*{Declarations}

\paragraph{Funding.} No funding was received for this work.
\paragraph{Competing interests.} The author declares no competing interests.
\paragraph{Author contributions.} The sole author conceived the method,
implemented the library and experiments, and wrote the manuscript.
\paragraph{Data and code availability.} All code, experiment scripts and
generated results are available in the project repository, to be disclosed
upon acceptance to preserve the double-anonymous review process; a private
link is available to editors and reviewers on request. The corpora referenced
in the manuscript's Related Work (CUHK Avenue, UCSD Ped2, TACO, ZeroWaste,
UAVVaste, AerialWaste, MARIDA, MADOS, Mivia-IWDD) are distributed by their
respective authors under their own terms; no data from them is redistributed
here.
\paragraph{Ethical approval.} Not applicable. No human-subjects data was
collected for this work. The real video corpora used (CUHK Avenue, UCSD Ped2)
are long-standing, publicly released video-anomaly-detection benchmarks
distributed by their original authors for research use; this paper is a
downstream, secondary use of that already-public release, not new data
collection, and involves no interaction with human subjects. This work
concerns the statistical calibration of alerting systems, not the governance
of their deployment; the manuscript's Limitations section gives a substantive
treatment of the human-oversight, accountability, and disclosure questions
that automated surveillance of public space raises, rather than treating the
topic as out of scope.

\end{document}
